\begin{abstract}
4D perception models, which classify 3D point clouds that evolve over time, are being built into robots, assistive monitoring and AR, yet before deploying one a practitioner cannot check \emph{which cue} it relies on: how a body part moves, or what shape it has. Existing explanations do not fit this setting. 3D attribution is single-frame, video attribution is confined to RGB pixels, gradient maps need white-box access, attention is absent from state-space backbones, and point-level perturbations assume a point correspondence across frames that real depth and LiDAR data lack. We present \methodlong (\method), a black-box, interventional attribution method that queries only the model's outputs. Its \frz operator removes a region's motion while keeping its shape, and a shape-free proxy removes the shape while keeping the motion. Together they span a motion$\times$geometry cube for every tracked region, whose order-symmetric Shapley decomposition is complete and whose interaction is measured rather than assumed away. The operators rewrite per-frame point \emph{sets}, so they need no correspondences, and a context-averaged estimator keeps the scores reliable when a structure spans several regions. We validate \method on the official PSTNet, P4Transformer and Mamba4D trained on MSR-Action3D. The decisive test is an audit: two binary tasks are built from real clips whose classes differ only in the motion or only in the shape of the same arm, and the motion share of the arm's importance follows the defining cue on all three backbones (\NumXMotShare\ vs.\ \NumXGeoShare\ on P4Transformer, \NumXPstMotShare\ vs.\ \NumXPstGeoShare\ on PSTNet, \NumXMbaMotShare\ vs.\ \NumXMbaGeoShare\ on Mamba4D), a distinction no single attribution map can make. The split costs no faithfulness: the necessity maps have the lowest deletion AUC on P4Transformer and PSTNet, tie with the best baseline on Mamba4D, and lose their correlation with the trained model's once the weights are randomised. Skeleton ground truth, re-sampling stability and a cost study complete the validation.
\end{abstract}
